% !TEX root = ../main.tex
\chapter{Хэрэгжилт}\label{ch:impl}

\section{Хэл, сан, орчин}
Java 17 (\code{javac --release 17}) дээр зөвхөн стандарт сан ашиглан бичсэн. Thread-д \code{java.lang.Thread},
Lockstep-д \code{java.util.concurrent.Semaphore}, web серверт JDK-д багтсан \code{com.sun.net.httpserver},
UI-д framework-гүй HTML/CSS/JavaScript (\code{<canvas>}) ашигласан. Maven, Gradle, JUnit хэрэглээгүй.
Build, ажиллуулах, тест, хэмжилтийг \code{build.sh}, \code{run.sh}, \code{test.sh}, \code{bench.sh} скриптүүд
(Windows: \code{.bat}) хийнэ.

\section{Parent, sorting ба merging thread}
Parent-ийн логик (Код~\ref{lst:parent}) номын тайлбарыг шууд дагадаг: хуваах, хоёр sorter-ийг эхлүүлэх,
хоёуланг нь хүлээх, merger-ийг эхлүүлэх, хүлээх.

\begin{lstlisting}[caption={Parent thread (\code{MultithreadedSorter.sort}, trace ба gate-ийн мөрийг товчлов)},label={lst:parent}]
int mid = (n + 1) / 2;                          // ceil(n/2)
SharedData data = new SharedData(input);        // list[] and result[]
SortingThread left  = new SortingThread(data, new SortParams("Sorter-L", 0, mid, la, leftAlgo), ...);
SortingThread right = new SortingThread(data, new SortParams("Sorter-R", mid, n, ra, rightAlgo), ...);
left.start();  right.start();
left.join();   right.join();
MergingThread merger = new MergingThread(data, new MergeParams("Merger", 0, mid, n), ...);
merger.start();
merger.join();
return new SortResult(input.clone(), data.result, mid, left, right, merger, sortPhase, total);
\end{lstlisting}

\code{SortingThread.run()} нь параметрээр авсан \code{[startIndex, endIndex)} мужийг сонгосон алгоритмаар
эрэмбэлнэ. \code{MergingThread} нь зүүн хэсгийн заагч \code{i}, баруун хэсгийн заагч \code{j}-ийг харьцуулж багыг
нь \code{result[k]}-д бичнэ (тэнцүү үед зүүнийхийг авна). Thread дотор алдаа гарвал тэр нь parent руу автоматаар
дамжихгүй тул thread бүр алдааг хадгалж, parent join хийсний дараа шалгадаг.

Алгоритмууд массивт шууд хандахын оронд \code{Ops} объектын \code{compare}, \code{swap}, \code{set} аргуудыг
дууддаг. Thread бүр өөрийн \code{Ops}-той тул үйлдлийн тоолуурт lock хэрэггүй. Trace асаалттай бол үйлдэл бүрийг
event болгон бүртгэж, Lockstep горимд үйлдэл бүрийн өмнө permit хүлээнэ. Хоёулаа унтраалттай үед энэ нь ердийн
массивын хандалт болно.

\section{Оролт ба event trace}
\code{InputParser} нь таслал, зай, шинэ мөрөөр тусгаарласан утга бүрийг бүхэл, бутархай, тоо биш гэж ангилж,
\code{int} мужийг шалгана. Алдааны мессеж байрлалыг нэрлэнэ (жишээ: “3-р утга 'abc' нь бүхэл тоо биш.”).
\code{Limits} нь \ref{sec:limits}-р хэсгийн хязгаарыг шалгаж, сондгой $n$-ийн үед тэмдэглэл гаргана. Оролтын
алдаанд CLI exit code 2, амжилттай үед 0 буцаана.

\code{TraceRecorder} нь event бүрт \code{AtomicLong}-оос дэс дугаар өгч, lock-гүй дараалалд нэмнэ. Event нь
\code{\{seq, t, th, type, a, b, c\}} бүтэцтэй (\code{CMP}, \code{SWP}, \code{SET}, \code{MCMP}, \code{MW},
\code{START}, \code{JOIN} гэх мэт) бөгөөд \code{--trace FILE}-аар JSON болгон хадгална. Хоосон талбарын
тэмдэглэгээг эхэндээ \code{Integer.MIN\_VALUE} гэж сонгосон нь бодит утгатай давхцаж байсан тул \code{long} талбар
ба \code{Long.MIN\_VALUE} болгож зассан.

\section{GitHub ба эх код}
Эх код \url{\RepoURL} хаягт байрладаг. Ажлыг үе шатад хувааж (core, алгоритм, оролт, trace,
web UI, Lockstep, тест, хэмжилт, баримт бичиг), үе шат бүрийн төгсгөлд зөвхөн хэрэгтэй файлыг сонгон Монгол
мессежтэй commit хийсэн. Эцсийн хувилбарыг \code{v1.0.1} tag бүхий GitHub Release болгон нийтэлсэн. Build-ийн файлуудыг \code{.gitignore}-оор хассан.
Repository-д \code{src/} (код), \code{tests/} (тест, хэмжилт, үр дүнгийн хүснэгт), \code{examples/} (тестийн
өгөгдөл), \code{docs/} (гарын авлага), \code{report/} (тайлан) хавтас байна.
